% Springer Nature manuscript prepared for the International Journal of
% Data Science and Analytics using the supplied sn-jnl class.
\documentclass[pdflatex,sn-mathphys-num]{sn-jnl}

\usepackage{graphicx}
\graphicspath{{../figures/final/}}
\usepackage{booktabs}
\usepackage{amsmath,amssymb,amsfonts}
\usepackage{microtype}
\usepackage{array}
\usepackage{url}
\setlength{\emergencystretch}{3em}
\raggedbottom

\begin{document}

\title[A Reproducible Event-Fixed Protocol for Evaluating Financial Stress Detectors]{A Reproducible Event-Fixed Protocol for Evaluating Financial Stress Detectors}

\author*[1]{\fnm{Prashant} \sur{Singh}}\email{prashants18488@gmail.com}
\author[2]{\fnm{Pranav} \sur{Singh}}

\affil[1]{\orgdiv{Department of Mathematics and Computing}, \orgname{Indian Institute of Technology Ropar}, \orgaddress{\city{Rupnagar}, \state{Punjab}, \country{India}}}
\affil[2]{\orgdiv{Department of Mathematics and Computing}, \orgname{Indian Institute of Technology Ropar}, \orgaddress{\city{Rupnagar}, \state{Punjab}, \country{India}}}

\abstract{Financial stress detection is important for risk monitoring, market surveillance, and decision support, but consistent evaluation is difficult because stress events are rare, event boundaries are ambiguous, and detector outputs are often compared under incompatible labeling rules. This paper presents a reproducible event-fixed evaluation protocol for daily financial stress detectors. The protocol fixes documented stress events before detector inspection, applies deterministic preprocessing, assigns rule-based drawdown severity categories, and reports pointwise, event-level, delay, uncertainty, multi-asset, and cost-sensitive metrics. As an illustrative use case, we compare a volatility percentile trigger with a hybrid retrospective confirmation rule that requires local instability evidence. On SPY, the hybrid reduces false positive rate (FPR) from 0.044 to 0.005 and increases precision from 0.640 to 0.886, but recall falls from 0.225 to 0.122 and F1 falls from 0.333 to 0.214. The same qualitative pattern appears across SPY, QQQ, IWM, HYG, and TLT: confirmation reduces false alarms but lowers recall and F1 and often delays first detection. The results show that event-fixed evaluation changes detector interpretation: the same confirmation rule can appear attractive under false-alarm-sensitive metrics while remaining weaker for recall, event coverage, and timeliness.}

\keywords{financial stress detection; event-fixed evaluation; time-series anomaly detection; financial risk monitoring; detection delay; cost-sensitive evaluation}

\maketitle

\section{Introduction}
Financial stress detectors are used in risk monitoring, market surveillance, financial decision support, and alert systems that must operate under asymmetric operational costs. False alarms can create alert fatigue, unnecessary escalation, and inefficient manual review, while missed stress episodes can delay risk mitigation or leave decision makers without timely warning. These competing costs make evaluation as important as model design: a detector that is attractive for conservative confirmation may be unsuitable for broad surveillance.

Consistent evaluation is difficult because financial stress events are rare and their boundaries are often ambiguous. Event dates may be documented after the fact, stress can begin before a named crisis date and persist afterward, and different studies may use different windows, labels, and output definitions. Data-dependent crisis selection, inconsistent event-window construction, false positives outside stress periods, detection delay, and precision--recall tradeoffs can all change the apparent ranking of detectors. The event-study literature emphasizes pre-specified dates for measuring market responses \cite{brown1985,mackinlay1997}, while time-series anomaly benchmarks show that label and scoring design can strongly affect conclusions \cite{lavin2015,wu2022}.

Many anomaly-detection, change-point, volatility, and critical-transition methods can be applied to financial time series. The evaluation layer has received less attention: which stress events define the target, how event windows are constructed, whether delay is measured explicitly, and how false alarms are weighed against missed episodes. For financial stress detection, this protocol layer is especially important because rare events and broad non-event periods can make a single headline metric misleading.

This paper develops an event-fixed evaluation protocol for financial stress detectors. The primary contribution is the protocol itself: a fixed panel of 13 documented U.S. equity-market stress episodes, deterministic adjusted-close preprocessing, descriptive realized-drawdown severity categories, pointwise and event-level metrics, first-detection delay and alignment measures, exploratory event and block bootstrap uncertainty, multi-asset validation, auxiliary diagnostic validation, and pointwise and event-aware cost-sensitive analysis. Calendar event dates are mapped to the nearest valid trading day for each asset, and detector thresholds are calibrated on each asset's own initial training segment without per-asset tuning.

The detector comparison serves as a case study of the protocol's interpretive value. A volatility-plus-instability hybrid sharply reduces false positives and improves precision relative to a volatility percentile trigger, but it lowers recall and F1. It also tends to delay alarms: the first in-window hybrid detection is often later than the corresponding volatility alarm, and event-level analysis shows that the hybrid detects fewer fixed events. Cost analysis clarifies the operational interpretation. Pointwise costs can favor the hybrid when non-event false positives dominate, but event-aware costs can favor volatility when missed events or late detection are penalized. The protocol therefore characterizes operating tradeoffs rather than declaring a universally superior detector.

\subsection*{Protocol contribution}
The protocol integrates established tools from event studies, volatility measurement, drawdown measurement, bootstrap uncertainty, and anomaly-detection evaluation. Its contribution is to combine pointwise metrics, event-level metrics, delay metrics, descriptive severity organization, multi-asset consistency checks, cost-sensitive interpretation, and a reproducibility package under pre-specified financial-stress event definitions. Table~\ref{tab:protocol_comparison} summarizes this integration.

\begin{table*}[t]
\centering
\caption{Relationship between established evaluation elements and the event-fixed protocol.}
\label{tab:protocol_comparison}
\scriptsize
\begin{tabular}{p{0.30\linewidth}p{0.62\linewidth}}
\toprule
Established evaluation element & Role in this protocol\\
\midrule
Fixed event panel & Uses a fixed panel of documented U.S. equity-market stress events before detector inspection.\\
Pre-specified inclusion rule & Applies explicit inclusion criteria based on broad market relevance, clear calendar dating, documented sources, and observable equity-market impact.\\
Event-level metrics & Reports whether each documented stress episode is detected, in addition to pointwise precision, recall, F1, and FPR.\\
Delay metrics & Reports first-detection delay and alignment error relative to the mapped event date.\\
Severity organization & Uses realized-drawdown terciles as descriptive categories for organizing event-level interpretation.\\
Multi-asset consistency check & Applies the same protocol without asset-specific detector tuning across equity, credit, and Treasury exchange-traded funds.\\
Cost-sensitive interpretation & Reports pointwise and event-aware operating costs to distinguish alert burden from missed-event and delay penalties.\\
Reproducibility package & Provides event definitions, source strings, mapped event tables, configuration, random seed, and reproduction procedure.\\
\bottomrule
\end{tabular}
\end{table*}

The contributions are:
\begin{enumerate}
\item an event-fixed financial stress evaluation protocol with pre-specified stress episodes and deterministic preprocessing;
\item descriptive realized-drawdown severity categories and event-level metrics for detection, coverage, delay, and alignment;
\item multi-asset validation across equity, credit, and Treasury exchange-traded funds;
\item exploratory bootstrap uncertainty and cost-sensitive analysis for deployment-facing interpretation;
\item an illustrative empirical characterization of the false-positive reduction, recall loss, and delay induced by confirmation-based stress filters.
\end{enumerate}

The remainder of the paper is organized as follows. Section 2 positions the protocol relative to event studies, anomaly benchmarks, financial stress measurement, critical-transition signals, change-point detection, and kernel methods. Section 3 defines the event-fixed protocol and the illustrative detectors. Section 4 reports the primary financial-stress evidence, followed by secondary sensitivity analyses and auxiliary diagnostic validation. Sections 5 and 6 discuss the operational interpretation and scope constraints, Section 7 summarizes reproducibility details, and Section 8 concludes.

\section{Related Work}
Event studies provide a finance precedent for evaluating market behavior around fixed dates. Brown and Warner's daily stock-return event-study design \cite{brown1985} and MacKinlay's survey of event-study methodology \cite{mackinlay1997} show why pre-specified event dates are central to credible empirical comparison: the event definition must be fixed before outcomes are interpreted. Financial stress episodes are also documented through institutional chronologies and post-event reports, including the National Bureau of Economic Research recession chronology \cite{nber2023}, the Financial Crisis Inquiry Commission report on the 2008 crisis \cite{fcic2011}, and the CFTC--SEC report on the May 2010 Flash Crash \cite{cftcsec2010}. These sources support auditable event selection, but event-study methods were not designed to evaluate continuous detector outputs, false alarms outside event windows, missed stress episodes, or detection delay.

Anomaly-detection benchmarks address the scoring problem from a different direction. NAB introduced labeled windows for streaming anomaly evaluation \cite{lavin2015}, and Ahmad et al. use that setting to study real-time unsupervised anomaly detection \cite{ahmad2017}. At the same time, Wu and Keogh show that benchmark labels and artifacts can create misleading impressions of progress \cite{wu2022}. This literature shows that labels, windows, and scoring rules can determine apparent progress. However, general anomaly benchmarks are not tailored to documented financial stress events, and they typically do not separate the operational roles of surveillance, confirmation, and escalation.

Financial stress measurement motivates the empirical signals used in the protocol. Drawdown summarizes realized loss severity \cite{magdon2004}, while realized volatility measures market variation \cite{andersen2003}. Heavy tails, stylized return facts, and volatility clustering are well documented in financial returns \cite{mandelbrot1963,cont2001,engle1982,bollerslev1986}. This literature explains why volatility, drawdown, and instability are meaningful inputs for stress detection. It does not, however, resolve the evaluation problem: a detector can respond to volatility or distributional change while still being poorly aligned with named stress events or unsuitable for a particular operating cost structure.

Critical-transition, change-point, and kernel methods motivate the detector families used as the illustrative case study. Critical-transition research links regime changes to rising variance, autocorrelation, and slower recovery \cite{scheffer2009,dakos2012}, while Boettiger and Hastings emphasize limits to detecting such signals in finite time series \cite{boettiger2012}. Change-point detection supplies sequential tools such as CUSUM \cite{page1954}, and kernel two-sample testing motivates MMD as a nonparametric distribution-comparison statistic \cite{gretton2012}. These methods provide plausible detector components, but method availability does not solve the comparison problem. Without a fixed event panel and deployment-relevant scoring, detector rankings can reflect evaluation choices rather than meaningful differences in financial stress monitoring.

Together, these literatures motivate the protocol developed below. Event studies motivate fixed dates, anomaly benchmarks motivate careful label and scoring design, financial stress research motivates volatility and drawdown-based measures, and critical-transition/change-point/kernel methods motivate the detector families. What remains unresolved is the integration of these ideas into a transparent financial-stress evaluation standard. The gap addressed here is therefore not the absence of detectors, but the absence of a compact, reproducible framework that makes stress-event definitions, event windows, severity organization, delay, and operating costs explicit before detector performance is interpreted.

\section{Methodology}
The methodology is organized around the evaluation protocol rather than around a particular detector. The protocol first fixes the financial stress events and trading-calendar mapping, then defines descriptive severity categories, detector outputs, event-window and event-level metrics, operating costs, and uncertainty summaries. The volatility and hybrid rules illustrate how the protocol separates false-alarm reduction from event coverage and timeliness.

\subsection{Data and preprocessing}
The primary analysis uses SPY, and the multi-asset validation uses SPY, QQQ (Invesco QQQ Trust), IWM (iShares Russell 2000 ETF), HYG (iShares iBoxx \$ High Yield Corporate Bond ETF), and TLT (iShares 20+ Year Treasury Bond ETF). For every asset, the adjusted close column is used. Let $P_t$ denote the adjusted close. The deterministic preprocessing is
\begin{align}
r_t &= \log(P_t/P_{t-1}),\\
\sigma_t &= \mathrm{sd}(r_{t-19:t}),\\
z_t &= (r_t-\mu_{t,50})/\sigma_{t,50},\\
x_t &= 0.5r_t + 0.3\sigma_t + 0.2z_t.
\end{align}
Fixed calendar events are mapped separately for each asset to the nearest valid trading day in that asset's calendar. The processed sample lengths after rolling-window preprocessing are reported in Appendix~\ref{app:implementation_details}.

\subsection{Inclusion rule for fixed events}
The event panel is fixed before detector evaluation, and the event list is not selected based on detector outputs. An event must satisfy all of the following criteria: broad market relevance to SPY/S\&P 500, a clear calendar date that can be mapped to a trading day, documentation by a regulator, central bank, NBER chronology, exchange, or major financial literature, and observable impact on U.S. equity markets. Event-source documentation is treated as evidence for inclusion rather than as author judgment about detector performance; the included sources are summarized in Table~\ref{tab:events} and retained in the reproducibility package.

The list is intended as a fixed stress-event panel rather than an exhaustive chronology of all U.S. market disruptions. Events are excluded when they lack a clear event date, have primarily localized or sector-specific effects without broad U.S. equity-market relevance, or overlap so closely with another included event that the event windows would not provide separable evaluation targets. The event panel draws on the NBER recession chronology for recession dating \cite{nber2023}, the Financial Crisis Inquiry Commission report for institutional documentation of the 2008 crisis \cite{fcic2011}, the CFTC--SEC report for the Flash Crash \cite{cftcsec2010}, and Khandani and Lo's analysis for the 2007 quant crisis \cite{khandani2007}. Other event-source strings are retained in the reproducibility materials rather than used as separate bibliographic claims.

\begin{table*}[t]
\centering
\caption{Fixed SPY stress-event panel with mapped trading dates, source labels, and realized-drawdown severity categories.}
\label{tab:events}
\scriptsize
\setlength{\tabcolsep}{3pt}
\begin{tabular}{>{\raggedright\arraybackslash}p{0.045\linewidth}>{\raggedright\arraybackslash}p{0.14\linewidth}>{\raggedright\arraybackslash}p{0.09\linewidth}>{\raggedright\arraybackslash}p{0.09\linewidth}>{\raggedright\arraybackslash}p{0.13\linewidth}>{\raggedright\arraybackslash}p{0.305\linewidth}>{\raggedright\arraybackslash}p{0.11\linewidth}}
\toprule
ID & Event & Calendar & Trading & Type & Source label & Severity \\
\midrule
E01 & Dot-com selloff & 2000-04-14 & 2000-04-14 & bear market & Fed History; NBER chronology & 0.114 high\\
E02 & Sep. 11 reopening & 2001-09-17 & 2001-09-17 & exogenous shock & SEC/NYSE reopening documentation & 0.182 critical\\
E03 & WorldCom selloff & 2002-07-24 & 2002-07-24 & accounting & SEC WorldCom litigation; Fed History & 0.196 critical\\
E04 & Quant crisis & 2007-08-09 & 2007-08-09 & liquidity & Khandani--Lo (2007); Fed accounts & 0.090 medium\\
E05 & Lehman crisis & 2008-09-29 & 2008-09-29 & systemic & FCIC report; NBER chronology & 0.344 critical\\
E06 & Flash crash & 2010-05-06 & 2010-05-06 & micro\-structure & CFTC--SEC Flash Crash report & 0.123 high\\
E07 & U.S. downgrade & 2011-08-08 & 2011-08-08 & sovereign credit & S\&P downgrade; Fed commentary & 0.166 critical\\
E08 & China devaluation & 2015-08-24 & 2015-08-24 & global growth & Fed/BIS market commentary & 0.112 medium\\
E09 & Volmageddon & 2018-02-05 & 2018-02-05 & volatility & Cboe volatility-product discussion; SEC/CFTC commentary & 0.101 medium\\
E10 & Q4 2018 selloff & 2018-12-24 & 2018-12-24 & liquidity & Federal Reserve Financial Stability Report & 0.156 high\\
E11 & COVID crash & 2020-03-16 & 2020-03-16 & pandemic & Fed FSR May 2020; NBER chronology & 0.337 critical\\
E12 & Inflation selloff & 2022-06-13 & 2022-06-13 & inflation & Federal Reserve Financial Stability Report & 0.122 high\\
E13 & Bank stress & 2023-03-13 & 2023-03-13 & banking & Fed SVB review; Fed FSR 2023 & 0.069 medium\\
\bottomrule
\end{tabular}
\end{table*}

\subsection{Severity organization}
All 13 events receive descriptive severity categories based on realized SPY drawdown. These categories organize interpretation within the fixed event panel; they are not used as statistical tests and do not imply population-level inference. For mapped event index $e$, severity is
\begin{equation}
\mathrm{severity}(e)=\left|\min_{t\in[e-20,\,e+20]}\left(\frac{P_t}{\max_{u\in[e-20,\,t]}P_u}-1\right)\right|.
\end{equation}
Critical, high, and medium categories are assigned by terciles of this severity score across the 13 events. No manual severity adjustment is applied. These SPY-derived categories remain the primary event-panel organization throughout the paper. For multi-asset robustness checks, asset-specific drawdown severity scores are also computed after mapping each event to the asset's own trading calendar, but those asset-specific categories do not replace the primary SPY event-panel categories.

\subsection{Event-window evaluation}
The main positive mask uses $[e-60,e+60]$ trading days. A $\pm60$-day window is used because financial stress often develops over weeks rather than on a single trading day, and stress can persist after the named calendar event. The event window is therefore an evaluation region for stress identification, not a claim about exact crisis boundaries or a one-day prediction target. This wide window can blur distinctions among advance warning, retrospective confirmation, and post-event detection, so delay and alignment metrics are reported separately. To check dependence on this choice, we also evaluate $\pm20$, $\pm30$, and $\pm45$ trading-day windows.

\subsection{Event-level evaluation}
Pointwise masks summarize alarm days but do not directly answer whether a fixed stress event was detected at all. We therefore compute event-level metrics for each ticker, event, and method using the same $[e-60,e+60]$ trading-day window. The indicator \texttt{detected\_event} equals one if at least one alarm occurs inside the window. The variable \texttt{first\_detection\_delay} is the first in-window alarm date minus the mapped event date, measured in trading days; it is missing if the event is not detected. The variable \texttt{closest\_detection\_delay} is the in-window alarm closest to the mapped event date minus the mapped event date. The variable \texttt{event\_coverage} is the fraction of event-window trading days flagged by the method. The indicators \texttt{pre\_event\_detection} and \texttt{post\_event\_detection} record whether any alarm occurs in $[e-60,e-1]$ and $[e,e+60]$, respectively. These quantities complement precision, recall, F1, and FPR by separating event coverage, missed events, and timeliness.

\subsection{Detector definitions}
All thresholds are calibrated on the first 70\% of each series and reused unchanged. For any scalar score $q_t$, training standardization uses $z_q(t)=(q_t-\mu_{q,\mathrm{train}})/\sigma_{q,\mathrm{train}}$. Unless otherwise stated, smoothing is a trailing five-point rolling mean, $\bar q_t=|\{j:0\leq j\leq4,t-j\geq0\}|^{-1}\sum_{j=0}^{4}q_{t-j}$.
\begin{itemize}
\item \textbf{Volatility percentile trigger:} compute $\sigma_t=\mathrm{sd}(r_{t-19:t})$, standardize it on the calibration period, smooth it with the five-point trailing mean, and set $v_t=1$ when the smoothed score exceeds its 90th calibration percentile.
\item \textbf{CUSUM:} use standardized returns $z_r(t)$ and a two-sided zero-reference cumulative sum, $C_t^+=\max(0,C_{t-1}^+ + z_r(t))$, $C_t^-=\min(0,C_{t-1}^- + z_r(t))$, and $C_t=\max(C_t^+,-C_t^-)$. The smoothed $C_t$ is thresholded at its 90th calibration percentile.
\item \textbf{MMD:} for each $t$, compare pre-window $x_{t-30:t-1}$ and post-window $x_{t:t+29}$ with an RBF kernel \cite{gretton2012}. The bandwidth is the median absolute pairwise pre/post distance, lower bounded by $10^{-4}$. Let $K(x_i,x_j)=\exp[-\|x_i-x_j\|^2/(2\gamma^2)]$, where $\gamma$ is the bandwidth; $\bar K_{aa}$, $\bar K_{bb}$, and $\bar K_{ab}$ denote the mean within-pre, within-post, and cross-window kernel values, respectively. The biased estimate $\widehat{\mathrm{MMD}}^2=\bar K_{aa}+\bar K_{bb}-2\bar K_{ab}$ is clipped at zero, smoothed with a trailing three-point mean, standardized on the calibration period, and then used in the instability score.
\item \textbf{AR(1) persistence statistic ($\kappa$):} over a trailing 36-day autoregressive window, estimate $\widehat\phi_t=(\sum a_i b_i)/(\sum a_i^2+10^{-10})$ using adjacent observations $a_i=x_i$, $b_i=x_{i+1}$. Define $\kappa_t=1-\widehat\phi_t$ and standardize $\kappa_t$ on the calibration period.
\item \textbf{Instability score:} $s(t)=0.5z_{\mathrm{MMD}}(t)+0.5|z_\kappa(t)|$. The smoothed score is thresholded at its 90th calibration percentile. MMD-only and $|\kappa|$-only confirmations use the same standardization and threshold rule for the corresponding component.
\item \textbf{Hybrid $W=10$:} let $v(t)$ be the binary volatility alarm and $i(t)$ be the binary instability alarm. The hybrid alarm is active only when $v(t)=1$ and at least one instability alarm occurs within $\pm W$ trading days:
\begin{equation}
h_W(t)=v(t)\wedge \max_{\tau\in[t-W,t+W]}i(\tau).
\end{equation}
Here, $i(\tau)$ denotes the binary instability alarm indicator obtained by thresholding the smoothed instability score $s(\tau)$ at its 90th calibration percentile. Because the default rule uses a symmetric $[t-W,t+W]$ confirmation window, it uses future information relative to date $t$ and should be interpreted as retrospective confirmation in the event-fixed evaluation, not as a causal real-time decision rule.
\end{itemize}
CIS is computed as a diagnostic trailing 10-day sum of negative $\kappa$ slopes and is included only in ablations. Appendix~\ref{app:implementation_details} summarizes implementation settings used to reproduce the analysis.

\subsection{Auxiliary validation analyses}
Synthetic and NAB analyses are included as auxiliary validation of qualitative detector behavior. The synthetic analysis checks whether the pipeline responds to controlled regime changes under noisy, non-Gaussian dynamics; it is not intended to establish real-market generalization. NAB provides an external qualitative validation setting outside SPY. Official NAB scores are not used because the paper compares detectors using the same pointwise precision, recall, F1, and FPR metrics used for the financial stress protocol. Implementation details for these auxiliary analyses are provided in Appendix~\ref{app:implementation_details}.

\subsection{Multi-asset validation}
The multi-asset validation applies the same preprocessing, detector definitions, event-window construction, and hyperparameters to QQQ, IWM, HYG, and TLT. Thresholds are calibrated on the first 70\% of each asset's own processed series, matching the SPY rule, but no detector parameter is tuned separately by asset. The purpose is to test whether the SPY false-positive/recall/delay tradeoff is specific to SPY or appears across related equity, credit, and Treasury proxies.

\subsection{Metrics, cost analysis, and bootstrap}
Pointwise precision, recall, F1, and FPR use the union of event windows. Event coverage is the fraction of event windows containing at least one detection. Event delay is signed first-detection delay relative to the event center, and event alignment error is the absolute distance to the nearest detection. If an event has no detection, it contributes zero coverage and missing delay/alignment; overlapping windows are merged only for pointwise metrics.

The pointwise cost analysis retains the rule
\begin{equation}
C_{\mathrm{point}}=c_{\mathrm{FP}}\mathrm{FP}+c_{\mathrm{FN}}\mathrm{FN},
\end{equation}
with $c_{\mathrm{FN}}=1$ and $c_{\mathrm{FP}}/c_{\mathrm{FN}}\in\{1,2,5,10,20,50\}$. Because this cost is computed over the pointwise event-window mask, it can be dominated by the many non-event days. We therefore also report an event-aware cost,
\begin{equation}
C_{\mathrm{event}}=c_{\mathrm{FA}}A_{\mathrm{out}}+c_{\mathrm{ME}}M+c_{\mathrm{D}}D_+,
\end{equation}
where $A_{\mathrm{out}}$ is the number of alarm days outside all event windows, $M$ is the number of missed events, and $D_+$ is the sum of $\max(\text{first detection delay},0)$ over detected events. Missed events are not given a separate delay penalty. We set $c_{\mathrm{ME}}=1$, evaluate $c_{\mathrm{FA}}/c_{\mathrm{ME}}\in\{0.01,0.05,0.1,0.25,0.5,1.0\}$, and examine $c_{\mathrm{D}}\in\{0,0.05,0.1\}$. These costs are illustrative deployment regimes, not estimates of universal economic utility.

We use 2000 bootstrap resamples and percentile 95\% confidence intervals as exploratory uncertainty characterization. The event bootstrap resamples the 13 event rows with replacement, rebuilds the event-window mask, and evaluates fixed detector outputs. The block bootstrap uses overlapping non-circular 20-day blocks; sampled blocks are concatenated until the original timeline length is reached, then truncated. Detector outputs are not recomputed in either bootstrap; uncertainty is over evaluation sampling rather than detector fitting. Bootstrap resampling cannot compensate for the small number of fixed events, so these intervals should be interpreted as descriptive uncertainty estimates rather than population-level generalization.

\section{Results}
The results are organized by evidentiary role. The primary evidence consists of the main SPY evaluation, multi-asset validation, event-level analysis, and cost-sensitive interpretation. Secondary sensitivity analyses then examine event-window width, descriptive severity organization, confirmation-window behavior, and exploratory bootstrap uncertainty. Synthetic and NAB analyses are reported last as auxiliary diagnostic validation rather than as additional evidence of financial-market generalization.

\subsection{Primary evidence: main SPY evaluation}
Table~\ref{tab:spy} gives the main $\pm60$ trading-day event-window evaluation. The hybrid has the lowest FPR and highest precision, but lower recall and F1 than the volatility trigger. Because the hybrid uses symmetric $\pm W$ confirmation, these results characterize retrospective confirmation under the fixed-event protocol rather than causal real-time alerting.
Figure~\ref{fig:spy_timeline} shows the fixed event windows and hybrid outputs on the SPY timeline.

\begin{table*}[t]
\centering
\caption{Main SPY event-fixed evaluation under the $\pm60$ trading-day event window.}
\label{tab:spy}
\begin{tabular}{lrrrrrrr}
\toprule
Method & Precision & Recall & F1 & FPR & Coverage & Delay & Align. err.\\
\midrule
Volatility percentile & 0.640 & 0.225 & 0.333 & 0.044 & 0.769 & -1.2 & 3.6\\
CUSUM & 0.308 & 0.445 & 0.364 & 0.344 & 0.462 & -60.0 & 0.2\\
Instability score & 0.433 & 0.168 & 0.242 & 0.076 & 0.923 & -25.2 & 20.3\\
MMD only & 0.411 & 0.164 & 0.235 & 0.081 & 0.846 & -16.2 & 16.8\\
$|\kappa|$ only & 0.292 & 0.110 & 0.160 & 0.092 & 0.769 & -26.2 & 20.5\\
Hybrid ($W=10$) & 0.886 & 0.122 & 0.214 & 0.005 & 0.615 & 6.8 & 8.3\\
\bottomrule
\end{tabular}
\end{table*}

Volatility fires close to the mapped event center on average, with mean delay $-1.2$ trading days. The hybrid fires later, with mean delay $+6.8$ trading days, indicating that confirmation reduces false alarms partly by waiting for additional instability evidence. Its lower FPR therefore comes at the cost of lower recall and less timely first detection. CUSUM attains high recall and F1, but its FPR of 0.344 and mean delay near $-60$ days indicate broad alarm activity near the start of the wide event windows rather than precise event-centered localization.

\begin{figure}[t]
\centering
\includegraphics[width=\linewidth]{timeline_spy.png}
\caption{SPY timeline with fixed event windows and hybrid outputs, showing selective retrospective confirmation within documented stress regions.}
\label{fig:spy_timeline}
\end{figure}

\subsection{Primary evidence: multi-asset validation}
Table~\ref{tab:multiasset_pointwise} and Figure~\ref{fig:multiasset_tradeoff} apply the same event-fixed protocol to all processed assets and compare the volatility trigger with the default hybrid. The SPY pattern appears in the same direction: the hybrid reduces FPR for every asset, but it also lowers recall and F1 for every asset. The magnitude varies. For QQQ, FPR falls from 0.039 to 0.006 while F1 falls from 0.280 to 0.139. For IWM, the FPR reduction is smaller, from 0.056 to 0.028, and F1 falls from 0.272 to 0.159. HYG and TLT show the same low-FPR/lower-F1 pattern, but their delays differ: HYG's hybrid mean delay is slightly earlier than volatility, while TLT remains delayed relative to volatility. These results indicate a consistent directional tradeoff across the evaluated assets rather than a universal improvement.

\begin{table*}[t]
\centering
\caption{Multi-asset event-fixed pointwise evaluation for the volatility trigger and hybrid confirmation rule.}
\label{tab:multiasset_pointwise}
\scriptsize
\begin{tabular}{llrrrrrr}
\toprule
Ticker & Method & Precision & Recall & F1 & FPR & Mean delay & Align. err.\\
\midrule
SPY & Volatility & 0.640 & 0.225 & 0.333 & 0.044 & -1.2 & 3.6\\
SPY & Hybrid & 0.886 & 0.122 & 0.214 & 0.005 & 6.8 & 8.3\\
QQQ & Volatility & 0.616 & 0.181 & 0.280 & 0.039 & -19.3 & 3.2\\
QQQ & Hybrid & 0.812 & 0.076 & 0.139 & 0.006 & 10.4 & 15.2\\
IWM & Volatility & 0.529 & 0.183 & 0.272 & 0.056 & -0.3 & 4.4\\
IWM & Hybrid & 0.536 & 0.093 & 0.159 & 0.028 & 14.2 & 16.2\\
HYG & Volatility & 0.501 & 0.163 & 0.246 & 0.064 & 0.2 & 2.2\\
HYG & Hybrid & 0.650 & 0.110 & 0.188 & 0.024 & -1.0 & 2.3\\
TLT & Volatility & 0.507 & 0.281 & 0.362 & 0.085 & -25.3 & 4.7\\
TLT & Hybrid & 0.367 & 0.094 & 0.150 & 0.051 & -4.8 & 23.5\\
\bottomrule
\end{tabular}
\end{table*}

\begin{figure}[t]
\centering
\includegraphics[width=\linewidth]{multi_asset_tradeoff.png}
\caption{Multi-asset FPR--F1 tradeoff under the common event-fixed protocol without asset-specific detector tuning.}
\label{fig:multiasset_tradeoff}
\end{figure}

\subsection{Primary evidence: event-level analysis}
Pointwise recall and FPR do not indicate whether individual stress episodes are detected at least once. Table~\ref{tab:spy_event_level} and Figure~\ref{fig:event_level_delay} report the SPY event-level comparison for volatility and hybrid. The hybrid detects 8 of 13 SPY events, compared with 10 of 13 for volatility. It misses the U.S. downgrade and inflation selloff events that volatility detects, while both methods miss the quant crisis, Volmageddon, and bank stress events under the fixed $\pm60$-day window. For detected events, hybrid coverage is usually lower, with equal or later first detection in several cases; the Lehman event is especially delayed for the hybrid, with first detection at +33 trading days.

\begin{table*}[t]
\centering
\caption{SPY event-level detection, first-detection delay, and event-window coverage for volatility and hybrid outputs.}
\label{tab:spy_event_level}
\scriptsize
\begin{tabular}{llrrrrrr}
\toprule
Event & Severity & Vol. det. & Hyb. det. & Vol. delay & Hyb. delay & Vol. cov. & Hyb. cov.\\
\midrule
Dot-com selloff & high & 1 & 1 & 4 & 4 & 0.241 & 0.096\\
September 11 reopening & critical & 1 & 1 & 6 & 6 & 0.132 & 0.124\\
WorldCom selloff & critical & 1 & 1 & -8 & -4 & 0.446 & 0.306\\
Quant crisis & medium & 0 & 0 & NA & NA & 0.000 & 0.000\\
Lehman crisis & critical & 1 & 1 & -7 & 33 & 0.562 & 0.231\\
Flash crash & high & 1 & 1 & 11 & 11 & 0.174 & 0.174\\
U.S. downgrade & critical & 1 & 0 & 2 & NA & 0.405 & 0.000\\
China devaluation & medium & 1 & 1 & 8 & 8 & 0.116 & 0.116\\
Volmageddon & medium & 0 & 0 & NA & NA & 0.000 & 0.000\\
Q4 2018 selloff & high & 1 & 1 & 4 & 4 & 0.124 & 0.124\\
COVID crash & critical & 1 & 1 & -8 & -8 & 0.405 & 0.405\\
Inflation selloff & high & 1 & 0 & -24 & NA & 0.331 & 0.000\\
Bank stress & medium & 0 & 0 & NA & NA & 0.000 & 0.000\\
\bottomrule
\end{tabular}
\end{table*}

The multi-asset event-level summary in Table~\ref{tab:multiasset_event_level} reinforces the same pattern. Hybrid event detection rates are lower than volatility for every asset: 0.615 versus 0.769 on SPY, 0.385 versus 0.462 on QQQ, 0.385 versus 0.538 on IWM, 0.231 versus 0.385 on HYG, and 0.308 versus 0.538 on TLT. Mean event-window coverage is also lower for the hybrid on every asset. Median first-detection delay is later on SPY, QQQ, IWM, and TLT; HYG is again the exception.

\begin{table}[t]
\centering
\caption{Multi-asset event-level detection rate, median delay, and mean event-window coverage.}
\label{tab:multiasset_event_level}
\scriptsize
\begin{tabular}{llrrr}
\toprule
Ticker & Method & Detect. rate & Med. delay & Mean cov.\\
\midrule
SPY & Vol. & 0.769 & 3.0 & 0.226\\
SPY & Hybrid & 0.615 & 5.0 & 0.121\\
QQQ & Vol. & 0.462 & -11.5 & 0.198\\
QQQ & Hybrid & 0.385 & 2.0 & 0.086\\
IWM & Vol. & 0.538 & 1.0 & 0.176\\
IWM & Hybrid & 0.385 & 5.0 & 0.090\\
HYG & Vol. & 0.385 & 2.0 & 0.123\\
HYG & Hybrid & 0.231 & -3.0 & 0.083\\
TLT & Vol. & 0.538 & -4.0 & 0.227\\
TLT & Hybrid & 0.308 & -2.5 & 0.076\\
\bottomrule
\end{tabular}
\end{table}

\begin{table}[t]
\centering
\caption{Severity-stratified event-level summary using descriptive realized-drawdown terciles from the SPY event panel.}
\label{tab:event_level_severity}
\scriptsize
\begin{tabular}{llrrr}
\toprule
Severity & Method & Detect. rate & Med. delay & Mean cov.\\
\midrule
Critical & Vol. & 0.840 & -2.6 & 0.317\\
Critical & Hybrid & 0.680 & 2.0 & 0.176\\
High & Vol. & 0.550 & -8.9 & 0.175\\
High & Hybrid & 0.300 & -1.3 & 0.066\\
Medium & Vol. & 0.150 & -26.0 & 0.046\\
Medium & Hybrid & 0.100 & -26.0 & 0.011\\
\bottomrule
\end{tabular}
\end{table}

\begin{figure*}[t]
\centering
\includegraphics[width=0.95\textwidth]{event_level_delay_coverage.png}
\caption{SPY event-level first-detection delay and event-window coverage, with missed events marked explicitly.}
\label{fig:event_level_delay}
\end{figure*}

\subsection{Primary evidence: cost-sensitive interpretation}
A simple pointwise cost model helps clarify why the precision/FPR tradeoff can matter operationally. With $c_{\mathrm{FN}}=1$ and $c_{\mathrm{FP}}/c_{\mathrm{FN}}\in\{1,2,5,10,20,50\}$, cost is $c_{\mathrm{FP}}\mathrm{FP}+c_{\mathrm{FN}}\mathrm{FN}$. Table~\ref{tab:cost} and Figure~\ref{fig:cost} show that the hybrid has lower SPY pointwise cost for all tested ratios, already marginally at 1:1 and increasingly as false positives become more expensive. Because the pointwise mask contains many non-event days, reducing false positives can dominate this cost even when recall, F1, event detection, and timeliness are worse.

\begin{table}[t]
\centering
\caption{Pointwise cost-sensitive comparison with $c_{\mathrm{FN}}=1$.}
\label{tab:cost}
\scriptsize
\begin{tabular}{rrrr}
\toprule
$c_{\mathrm{FP}}/c_{\mathrm{FN}}$ & Vol. cost & Hybrid cost & Preferred\\
\midrule
1 & 1384 & 1372 & Hybrid\\
2 & 1579 & 1396 & Hybrid\\
5 & 2164 & 1468 & Hybrid\\
10 & 3139 & 1588 & Hybrid\\
20 & 5089 & 1828 & Hybrid\\
50 & 10939 & 2548 & Hybrid\\
\bottomrule
\end{tabular}
\end{table}

The event-aware cost in Table~\ref{tab:event_cost} changes the interpretation. With $c_{\mathrm{D}}=0.05$, volatility is preferred for SPY when false-alarm days are cheap relative to missed events ($c_{\mathrm{FA}}/c_{\mathrm{ME}}=0.01$), because the hybrid misses more events and has larger positive delay. As false alarms become more costly, the hybrid becomes preferred. Thus, the cost analysis does not establish universal hybrid superiority. It shows that the preferred detector depends on whether the operating environment penalizes false-alarm days more heavily than missed or delayed stress events. The hybrid is more attractive when alert burden dominates, whereas the volatility trigger remains competitive when missed events and late confirmation are costly.

\begin{table}[t]
\centering
\caption{SPY event-aware cost comparison with $c_{\mathrm{ME}}=1$ and $c_{\mathrm{D}}=0.05$.}
\label{tab:event_cost}
\scriptsize
\begin{tabular}{lrrrr}
\toprule
$c_{\mathrm{FA}}/c_{\mathrm{ME}}$ & Vol. cost & Hybrid cost & Vol. missed & Hyb. missed\\
\midrule
0.01 & 6.70 & 8.54 & 3 & 5\\
0.05 & 14.50 & 9.50 & 3 & 5\\
0.10 & 24.25 & 10.70 & 3 & 5\\
0.25 & 53.50 & 14.30 & 3 & 5\\
0.50 & 102.25 & 20.30 & 3 & 5\\
1.00 & 199.75 & 32.30 & 3 & 5\\
\bottomrule
\end{tabular}
\end{table}

Across assets, the same cost-regime dependence appears: under event-aware cost with $c_{\mathrm{D}}=0.05$, volatility is preferred in at least one tested false-alarm/missed-event regime for every processed asset. This is consistent with the event-level findings: when missed events and delayed confirmation are penalized, a broader volatility trigger can be preferable despite its higher false-alarm rate.

\begin{figure*}[t]
\centering
\includegraphics[width=0.95\textwidth]{cost_comparison.png}
\caption{SPY cost comparison under pointwise and event-aware operating objectives.}
\label{fig:cost}
\end{figure*}

\subsection{Secondary robustness: event-window sensitivity}
Table~\ref{tab:eventwindow} and Figure~\ref{fig:eventwindow} show that the main tradeoff persists across event-window widths: the hybrid has lower FPR and higher precision, but lower recall and F1 than the volatility trigger for all tested widths. This analysis checks whether the interpretation depends on the default $\pm60$-day event window.

\begin{table}[t]
\centering
\caption{Event-window sensitivity for pointwise precision, recall, F1, and FPR.}
\label{tab:eventwindow}
\scriptsize
\begin{tabular}{llrrrr}
\toprule
Window & Method & Prec. & Recall & F1 & FPR\\
\midrule
$\pm20$ & Vol. & 0.390 & 0.396 & 0.393 & 0.061\\
$\pm20$ & Hybrid & 0.545 & 0.216 & 0.309 & 0.018\\
$\pm30$ & Vol. & 0.499 & 0.344 & 0.407 & 0.052\\
$\pm30$ & Hybrid & 0.692 & 0.186 & 0.293 & 0.012\\
$\pm45$ & Vol. & 0.571 & 0.266 & 0.363 & 0.048\\
$\pm45$ & Hybrid & 0.815 & 0.148 & 0.251 & 0.008\\
$\pm60$ & Vol. & 0.640 & 0.225 & 0.333 & 0.044\\
$\pm60$ & Hybrid & 0.886 & 0.122 & 0.214 & 0.005\\
\bottomrule
\end{tabular}
\end{table}

\begin{figure}[t]
\centering
\includegraphics[width=\linewidth]{event_window_sensitivity.png}
\caption{Event-window sensitivity across $\pm20$, $\pm30$, $\pm45$, and $\pm60$ trading-day windows.}
\label{fig:eventwindow}
\end{figure}

\subsection{Secondary robustness: severity, component, and confirmation-window analyses}
The descriptive severity categories organize the event-level interpretation within this fixed panel. Within this event panel, events assigned to the highest realized-drawdown tercile have F1 0.412 for volatility and 0.316 for hybrid, while FPR falls from 0.057 to 0.015. For events assigned to the middle tercile, F1 falls from 0.195 to 0.134 and FPR falls from 0.080 to 0.030. Events assigned to the lowest tercile have very low F1 for both methods. These tercile summaries are descriptive and should not be interpreted as statistical inference about severity effects.

The MMD-only value in Table~\ref{tab:spy} is the standalone thresholded MMD component. In contrast, the component ablation applies MMD-only as the confirmation signal inside the hybrid rule; that hybrid confirmation variant gives F1 0.219 versus 0.214 for the default no-CIS hybrid. The no-CIS score is more conservative because FPR is 0.005 versus 0.011. Figure~\ref{fig:window_sensitivity} shows that $W=10$ is not F1-optimal: F1 increases to 0.291 at $W=60$ while FPR also increases to 0.031. Because this analysis uses symmetric confirmation windows, it should be interpreted as retrospective sensitivity analysis.

\begin{figure}[t]
\centering
\includegraphics[width=\linewidth]{confirmation_window_sensitivity.png}
\caption{Hybrid confirmation-window sensitivity showing how the retrospective confirmation window changes F1 and FPR.}
\label{fig:window_sensitivity}
\end{figure}

\subsection{Secondary robustness: bootstrap uncertainty}
The event bootstrap gives hybrid-minus-volatility $\Delta$F1 $=-0.091$ with 95\% CI $[-0.166,0.005]$ and $\Delta$FPR $=-0.045$ with CI $[-0.058,-0.035]$. The block bootstrap gives $\Delta$F1 $=-0.120$ with CI $[-0.203,-0.047]$ and $\Delta$FPR $=-0.039$ with CI $[-0.062,-0.019]$. These bootstrap results provide exploratory uncertainty characterization: the hybrid's FPR difference is consistently negative in these resamples, whereas the F1 difference is negative or near zero. Figure~\ref{fig:bootstrap_violin} visualizes these bootstrap distributions.
The bootstrap analysis is kept on the primary SPY protocol. Because only 13 fixed events are available, bootstrap resampling cannot compensate for the limited event count; the intervals should be interpreted as descriptive uncertainty estimates rather than evidence of broad generalization. The present multi-asset analysis is interpreted as directional evidence of the qualitative tradeoff rather than as a pooled inferential estimate.

\begin{figure}[t]
\centering
\includegraphics[width=\linewidth]{bootstrap_delta.png}
\caption{Bootstrap distributions of hybrid-minus-volatility differences under the fixed-event protocol.}
\label{fig:bootstrap_violin}
\end{figure}

\subsection{Auxiliary diagnostic validation: synthetic and NAB analyses}
Figure~\ref{fig:synthetic_example} shows one generated sequence with latent labels and detector outputs.

\begin{figure*}[t]
\centering
\includegraphics[width=0.95\textwidth]{synthetic_example.png}
\caption{Synthetic diagnostic validation under controlled regime changes and noisy dynamics.}
\label{fig:synthetic_example}
\end{figure*}

The synthetic example contains stable, drift, and collapse regimes with heavy-tailed noise and volatility amplification near collapse. The plotted detector outputs show that the pipeline can respond to controlled regime changes under noisy, non-Gaussian dynamics. This diagnostic check establishes behavioral responsiveness in the illustrative generator, not real-market generalization.

On NAB, mean F1 is similar for hybrid and volatility (0.141 versus 0.140), while hybrid has lower FPR (0.105 versus 0.141). This resembles the qualitative pattern observed in SPY, where confirmation reduces false positives without producing a clear F1 improvement. Figure~\ref{fig:nab_results} visualizes this auxiliary validation by subset. Because official NAB scores are not used, the task is not finance-specific, and the NAB path uses a rolling mean-difference proxy for MMD, these results are interpreted as diagnostic consistency checks rather than leaderboard results or a strict detector-equivalence benchmark.

\begin{figure*}[t]
\centering
\includegraphics[width=0.95\textwidth]{nab_subset_results.png}
\caption{NAB external qualitative validation by subset using the same pointwise metrics as the financial-stress protocol.}
\label{fig:nab_results}
\end{figure*}

\begin{table}[t]
\centering
\caption{NAB subset-level F1 and FPR for volatility and hybrid outputs.}
\scriptsize
\setlength{\tabcolsep}{2pt}
\begin{tabular}{lrrrr}
\toprule
Subset & Vol. F1 & Vol. FPR & Hyb. F1 & Hyb. FPR\\
\midrule
artificialNoAnomaly & 0.000 & 0.095 & 0.000 & 0.040\\
artificialWithAnomaly & 0.168 & 0.107 & 0.163 & 0.080\\
realAWSCloudwatch & 0.108 & 0.166 & 0.108 & 0.119\\
realAdExchange & 0.148 & 0.191 & 0.151 & 0.169\\
realKnownCause & 0.179 & 0.232 & 0.194 & 0.185\\
realTraffic & 0.300 & 0.082 & 0.302 & 0.056\\
realTweets & 0.106 & 0.092 & 0.097 & 0.068\\
\bottomrule
\end{tabular}
\end{table}

\section{Discussion}
The results illustrate why financial stress detectors should be evaluated as operating regimes rather than ranked by a single aggregate metric. Under the event-fixed protocol, the hybrid occupies a different point on the surveillance-confirmation spectrum than the volatility trigger. The protocol makes this distinction visible by evaluating false alarms, event coverage, timing, and operating costs within the same fixed event structure.

For deployment, this distinction matters more than a marginal difference in any single score. A broad surveillance layer may tolerate additional alerts if doing so improves event coverage and timeliness. A confirmation layer may instead sacrifice coverage when escalation is costly and false positives impose review or governance burdens. The evaluated hybrid is therefore most naturally interpreted as a second-stage retrospective confirmation mechanism, whereas the volatility trigger is closer to a first-stage surveillance rule.

Event-fixed evaluation also changes the interpretation of delay. A detector can appear useful in pointwise scoring while still firing too late for prospective monitoring applications, or it can identify only a small fraction of a stress window while maintaining high precision. By reporting first-detection delay, alignment error, and event coverage, the protocol separates retrospective confirmation from broader stress surveillance.

The event-level and descriptive severity analyses show why the unit of evaluation matters. For crisis monitoring, the relevant question is not only how many individual days are classified correctly, but whether a documented stress episode is detected, when it is detected, and how much of the episode is covered. These measures bring evaluation closer to the operational risk-management problem than pointwise scores alone.

The multi-asset analysis reinforces the value of protocol-based evaluation. Stress is expressed differently across equity, credit, and Treasury proxies, so asset-level differences in delay and coverage are expected. The key contribution is that the same protocol exposes the same operating question across assets: whether a detector reduces alert burden by accepting weaker coverage or preserves broader surveillance at the cost of more false alarms.

The cost-sensitive results make the assumed deployment objective explicit. Pointwise costs emphasize alert burden over many non-event days, whereas event-aware costs penalize missed stress episodes and delayed confirmation directly. Reporting both views prevents the detector comparison from collapsing into a universal ranking and instead links detector choice to the institution's loss structure.

More broadly, the paper argues for event-fixed evaluation as a reporting standard for financial stress detection. Fixed event definitions reduce the risk of data-dependent crisis selection; documented sources improve auditability; descriptive severity categories support stratified interpretation within the event panel; and delay/alignment metrics clarify whether a detector is useful before, during, or after a stress episode. These protocol choices make detector comparisons more transparent and performance claims easier to interpret across deployment settings.

These implications are subject to the scope of the event panel, asset set, auxiliary validation, and uncertainty analysis. The next section states those constraints explicitly so that the protocol's conclusions are interpreted as transparent evidence from a fixed design rather than as universal detector claims.

\section{Limitations}
The fixed event panel is documented and reproducible, but it is not exhaustive. It focuses on 13 U.S. market stress episodes and does not cover all possible forms of financial stress, international crises, sector-specific disruptions, liquidity events, or credit-market episodes. A broader panel could change event-level detection rates and descriptive severity-stratified conclusions.

The primary severity categories are SPY-centered and exploratory. They are based on realized SPY drawdown around each mapped event date and therefore measure realized market impact rather than ex-ante severity or narrative importance. Because the categories are terciles within a 13-event panel, they are descriptive organizational labels rather than inferential severity classes. Asset-specific drawdown scores are computed for robustness, but the primary event categories may not perfectly match QQQ, IWM, HYG, or TLT, especially for assets with later inception dates or different economic exposures.

The asset set is broader than a single SPY experiment but remains limited. SPY, QQQ, and IWM cover major U.S. equity proxies, while HYG and TLT add credit and Treasury exposures, but the analysis does not include individual equities, international indexes, commodities, currencies, options-implied measures, or macro-financial indicators. No claim is made that the evaluated detector is state of the art; the goal is to evaluate detector behavior under a transparent protocol.

The synthetic data are illustrative rather than calibrated to real market dynamics. They are useful for diagnostic validation of pipeline responsiveness to controlled regime changes, but they do not establish real-world generalization. The NAB analysis is also qualitative: the benchmark is not finance-specific, official NAB scoring is not used, and the NAB implementation uses a faster proxy for MMD rather than the exact SPY computation.

The uncertainty and cost analyses are also limited. The bootstrap holds detector outputs fixed and quantifies evaluation-sampling uncertainty on the primary SPY protocol; it does not capture full model-selection uncertainty, hyperparameter-search uncertainty, or pooled multi-asset uncertainty, and it cannot compensate for the small number of fixed events. The event-aware cost model is deliberately simple and uses stylized penalties for false-alarm days, missed events, and positive delay. Real deployments would require institution-specific loss functions, escalation costs, and decision horizons.

The default hybrid rule uses a symmetric $\pm W$ confirmation window and therefore uses future information relative to date $t$. It should be interpreted as retrospective confirmation in the event-fixed protocol, not as a causal real-time decision rule. A causal trailing-only confirmation rule would be a useful future extension for deployment settings that require strictly prospective signals.

\section{Reproducibility and Code Availability}
The study is accompanied by a reproducibility package intended to support independent inspection and regeneration of the reported analyses. The package includes the fixed event panel, the per-asset mapped event table, full event-source strings, detector settings, preprocessing windows, bootstrap settings, and the random seed. The reported experiments use seed 42 and calibrate thresholds on the first 70\% of each processed asset series.

The reproducibility package is publicly available at \url{https://github.com/prashant290605/Financial-stress-event-fixed-protocol}. The repository includes package requirements, data-download instructions, expected runtimes, generated-output locations, and known reproduction limitations.

\section{Conclusion}
This paper presents a reproducible event-fixed evaluation protocol for financial stress detectors. By fixing the stress-event panel before detector inspection, mapping events to trading calendars, assigning descriptive severity categories, and reporting pointwise, event-level, delay, alignment, exploratory bootstrap, and cost-sensitive metrics, the protocol makes detector comparisons less dependent on a single aggregate score.

The empirical findings show how this protocol changes interpretation. A volatility-plus-instability retrospective confirmation rule can look attractive when false alarms are costly, but less attractive when recall, event coverage, and timeliness are emphasized. These results support interpreting the hybrid as a conservative confirmation or escalation layer within the post hoc event-fixed evaluation.

The central implication is that financial stress detectors should be evaluated against transparent event definitions and deployment-relevant metrics. In this study, the same detector can appear attractive under false-alarm-sensitive pointwise costs and less attractive when missed events or delayed detection are emphasized. The evaluation protocol therefore clarifies not only which detector performs better under a given metric, but also what operational role that metric implies.

\backmatter

\section*{Statements and Declarations}

\bmhead{Author Contributions}
Prashant Singh conceived the study, implemented the methodology, performed the experiments, analyzed the results, and prepared the original manuscript draft. Pranav Singh contributed to methodology refinement, interpretation of results, manuscript review, and supervision of the research process. Both authors reviewed and approved the final manuscript.

\bmhead{Funding}
The authors received no specific funding for this work.

\bmhead{Data Availability}
The financial market series used in this study are publicly available adjusted-close series for SPY, QQQ, IWM, HYG, and TLT. The financial-stress event records are documented by the sources summarized in Table~\ref{tab:events} and retained with the event panel in the reproducibility materials. The analysis can be reproduced using the accompanying code, event definitions, and data-download instructions.

\bmhead{Code Availability}
Code and reproducibility materials are publicly available at \url{https://github.com/prashant290605/Financial-stress-event-fixed-protocol}.

\bmhead{Competing Interests}
The authors declare no competing interests.

\clearpage
\begin{appendices}
\renewcommand{\theHtable}{A.\arabic{table}}
\renewcommand{\theHfigure}{A.\arabic{figure}}
\renewcommand{\theHequation}{A.\arabic{equation}}
\section{Implementation and Auxiliary Validation Details}
\label{app:implementation_details}
This appendix reports implementation details that support reproducibility and auditing. The fixed event panel is contained in \texttt{data/events\_spy.csv}, and the per-asset mapped event table is contained in \texttt{data/events\_multi\_asset.csv}. Full event-source strings are retained with the event panel. Detector settings, preprocessing windows, bootstrap settings, and the random seed are specified in \texttt{config/paper\_experiment.json}. The primary reproduction procedure is provided in \texttt{scripts/reproduce\_paper.py}; additional documented procedures are described in \texttt{SUPPLEMENTARY\_README.md}.

\begin{table}[t]
\centering
\caption{Multi-asset data availability after deterministic preprocessing.}
\label{tab:asset_availability}
\scriptsize
\begin{tabular}{llrr}
\toprule
Ticker & Asset & Dates & Days\\
\midrule
SPY & S\&P 500 ETF & 2000-03-15--2023-12-29 & 5987\\
QQQ & Nasdaq 100 ETF & 2000-03-15--2023-12-29 & 5987\\
IWM & Russell 2000 ETF & 2000-08-08--2023-12-29 & 5886\\
HYG & High-yield bond ETF & 2007-06-21--2023-12-29 & 4161\\
TLT & Treasury bond ETF & 2002-10-09--2023-12-29 & 5343\\
\bottomrule
\end{tabular}
\end{table}

\begin{table*}[t]
\centering
\caption{Detector and evaluation hyperparameters used in the reported analyses.}
\label{tab:hyperparams}
\scriptsize
\begin{tabular}{p{0.18\linewidth}p{0.76\linewidth}}
\toprule
Component & Setting \\
\midrule
Asset preprocessing & Adjusted close; log returns; 20-day rolling volatility; 50-day return z-score; scalar signal $x_t=0.5r_t+0.3\sigma_t+0.2z_t$.\\
Calibration & First 70\% of each series; training-period z-scores; 90th percentile alarm thresholds.\\
Volatility trigger & 20-day volatility window; training z-score; five-point trailing mean; 90th percentile threshold.\\
CUSUM & Input: standardized returns; two-sided zero-reference recursion with no separate drift/reference parameter; five-point trailing mean; 90th percentile threshold.\\
MMD & Pre/post windows of 30 observations; radial basis function (RBF) kernel; median pairwise distance bandwidth with $10^{-4}$ floor; training-period standardization.\\
AR(1) persistence statistic ($\kappa$) & Trailing 36-day AR(1) window; $\widehat\phi_t=(\sum a_i b_i)/(\sum a_i^2+10^{-10})$; $\kappa_t=1-\widehat\phi_t$; training-period standardization.\\
Hybrid & Default $W=10$ trading days; alarm only when volatility is active and instability appears within $[t-W,t+W]$.\\
Bootstrap & 2000 resamples; event bootstrap over event rows; block bootstrap with overlapping non-circular 20-day blocks; percentile 95\% confidence intervals.\\
\bottomrule
\end{tabular}
\end{table*}

Synthetic data are used only for controlled diagnostic validation; SPY and NAB thresholds are calibrated on the first 70\% of each respective series, and no synthetic-trained weights are transferred. Synthetic labels are generated from latent regimes: stable before stress onset, drift between stress onset and collapse onset, and collapse afterward. The signal follows
\begin{equation}
x_t=0.7x_{t-1}+d_t-\beta x_{t-1}^3+\epsilon_t,\quad
\epsilon_t=\sigma_t u_t+\ell_t,
\end{equation}
where $u_t$ is Student-$t$, $\ell_t$ is Laplace noise, $d_t<0$ is localized drift during the drift regime, and $\sigma_t$ is amplified in stress and near collapse.

\begin{table}[t]
\centering
\caption{Synthetic diagnostic generator settings.}
\label{tab:synthetic_params}
\scriptsize
\begin{tabular}{lr}
\toprule
Parameter & Value/range\\
\midrule
Series & 30\\
Length & 800\\
Student-$t$ df & [3, 8]\\
Laplace scale & [0.01, 0.05]\\
Calm volatility & [0.005, 0.02]\\
Stress multiplier & [1.5, 4.0]\\
Drift magnitude & [-0.08, -0.02]\\
Nonlinear $\beta$ & [0.01, 0.08]\\
\bottomrule
\end{tabular}
\end{table}

For the NAB auxiliary validation, all 58 series in the official \texttt{combined\_windows.json} are evaluated; anomaly windows are converted to binary masks, parameters are reused without tuning, and results are not used for leaderboard comparison. For runtime tractability on all 58 series, the NAB path uses a deterministic rolling mean-difference proxy for MMD rather than the quadratic exact-MMD implementation used for SPY. Because this proxy differs from exact MMD, NAB results are interpreted only qualitatively and not as a strict detector-equivalence benchmark.

\end{appendices}

\begin{thebibliography}{20}
\bibitem{brown1985} S.~J. Brown and J.~B. Warner, ``Using daily stock returns: The case of event studies,'' \emph{Journal of Financial Economics}, vol. 14, no. 1, pp. 3--31, 1985.
\bibitem{mackinlay1997} A.~C. MacKinlay, ``Event studies in economics and finance,'' \emph{Journal of Economic Literature}, vol. 35, no. 1, pp. 13--39, 1997.
\bibitem{lavin2015} A. Lavin and S. Ahmad, ``Evaluating real-time anomaly detection algorithms---the Numenta Anomaly Benchmark,'' in \emph{IEEE International Conference on Machine Learning and Applications}, 2015.
\bibitem{wu2022} R. Wu and E. Keogh, ``Current time series anomaly detection benchmarks are flawed and are creating the illusion of progress,'' in \emph{IEEE International Conference on Big Data}, 2022.
\bibitem{nber2023} National Bureau of Economic Research, ``U.S. business cycle expansions and contractions,'' 2023.
\bibitem{fcic2011} Financial Crisis Inquiry Commission, \emph{The Financial Crisis Inquiry Report}. U.S. Government Printing Office, 2011.
\bibitem{cftcsec2010} Commodity Futures Trading Commission and Securities and Exchange Commission, \emph{Findings Regarding the Market Events of May 6, 2010}, 2010.
\bibitem{magdon2004} M. Magdon-Ismail, A.~F. Atiya, A. Pratap, and Y.~S. Abu-Mostafa, ``On the maximum drawdown of a Brownian motion,'' \emph{Journal of Applied Probability}, vol. 41, no. 1, pp. 147--161, 2004.
\bibitem{andersen2003} T.~G. Andersen, T. Bollerslev, F.~X. Diebold, and P. Labys, ``Modeling and forecasting realized volatility,'' \emph{Econometrica}, vol. 71, no. 2, pp. 579--625, 2003.
\bibitem{scheffer2009} M. Scheffer et al., ``Early-warning signals for critical transitions,'' \emph{Nature}, vol. 461, pp. 53--59, 2009.
\bibitem{dakos2012} V. Dakos et al., ``Methods for detecting early warnings of critical transitions in time series illustrated using simulated ecological data,'' \emph{PLoS ONE}, vol. 7, no. 7, e41010, 2012.
\bibitem{boettiger2012} C. Boettiger and A. Hastings, ``Quantifying limits to detection of early warning for critical transitions,'' \emph{Journal of the Royal Society Interface}, vol. 9, pp. 2527--2539, 2012.
\bibitem{page1954} E.~S. Page, ``Continuous inspection schemes,'' \emph{Biometrika}, vol. 41, no. 1/2, pp. 100--115, 1954.
\bibitem{gretton2012} A. Gretton et al., ``A kernel two-sample test,'' \emph{Journal of Machine Learning Research}, vol. 13, pp. 723--773, 2012.
\bibitem{mandelbrot1963} B. Mandelbrot, ``The variation of certain speculative prices,'' \emph{Journal of Business}, vol. 36, no. 4, pp. 394--419, 1963.
\bibitem{cont2001} R. Cont, ``Empirical properties of asset returns: Stylized facts and statistical issues,'' \emph{Quantitative Finance}, vol. 1, no. 2, pp. 223--236, 2001.
\bibitem{engle1982} R.~F. Engle, ``Autoregressive conditional heteroscedasticity with estimates of the variance of United Kingdom inflation,'' \emph{Econometrica}, vol. 50, no. 4, pp. 987--1007, 1982.
\bibitem{bollerslev1986} T. Bollerslev, ``Generalized autoregressive conditional heteroskedasticity,'' \emph{Journal of Econometrics}, vol. 31, no. 3, pp. 307--327, 1986.
\bibitem{ahmad2017} S. Ahmad, A. Lavin, S. Purdy, and Z. Agha, ``Unsupervised real-time anomaly detection for streaming data,'' \emph{Neurocomputing}, vol. 262, pp. 134--147, 2017.
\bibitem{khandani2007} A.~E. Khandani and A.~W. Lo, ``What happened to the quants in August 2007?'' \emph{Journal of Investment Management}, vol. 5, no. 4, pp. 5--54, 2007.
\end{thebibliography}
\end{document}
